\documentclass[10pt,a4paper]{amsart}
\usepackage{amsmath,amssymb,amsthm,graphicx}
\usepackage{verbatim}
\newenvironment{reference}{\comment}{\endcomment}
%\newenvironment{reference}{}{}
\newenvironment{slogan}{\comment}{\endcomment}
\newenvironment{history}{\comment}{\endcomment}

% For commutative diagrams we use Xy-pic
\usepackage[all]{xy}

% We use 2cell for 2-commutative diagrams.
\xyoption{2cell}
% Theorem environments.
%
\theoremstyle{plain}
\newtheorem{theorem}[subsection]{Theorem}
\newtheorem{proposition}[subsection]{Proposition}
\newtheorem{lemma}[subsection]{Lemma}

\theoremstyle{definition}
\newtheorem{definition}[subsection]{Definition}
\newtheorem{example}[subsection]{Example}
\newtheorem{exercise}[subsection]{Exercise}
\newtheorem{situation}[subsection]{Situation}

\theoremstyle{remark}
\newtheorem{remark}[subsection]{Remark}
\newtheorem{remarks}[subsection]{Remarks}

\numberwithin{equation}{subsection}

% Macros
%
\def\lim{\mathop{\mathrm{lim}}\nolimits}
\def\colim{\mathop{\mathrm{colim}}\nolimits}
\def\Spec{\mathop{\mathrm{Spec}}}
\def\Hom{\mathop{\mathrm{Hom}}\nolimits}
\def\Ext{\mathop{\mathrm{Ext}}\nolimits}
\def\SheafHom{\mathop{\mathcal{H}\!\mathit{om}}\nolimits}
\def\SheafExt{\mathop{\mathcal{E}\!\mathit{xt}}\nolimits}
\def\Sch{\mathit{Sch}}
\def\Mor{\mathop{\mathrm{Mor}}\nolimits}
\def\Ob{\mathop{\mathrm{Ob}}\nolimits}
\def\Sh{\mathop{\mathit{Sh}}\nolimits}
\def\NL{\mathop{N\!L}\nolimits}
\def\CH{\mathop{\mathrm{CH}}\nolimits}
\def\proetale{{pro\text{-}\acute{e}tale}}
\def\etale{{\acute{e}tale}}
\def\QCoh{\mathit{QCoh}}
\def\Ker{\mathop{\mathrm{Ker}}}
\def\Im{\mathop{\mathrm{Im}}}
\def\Coker{\mathop{\mathrm{Coker}}}
\def\Coim{\mathop{\mathrm{Coim}}}

% Boxtimes
%
\DeclareMathSymbol{\boxtimes}{\mathbin}{AMSa}{"02}

%
% Macros for moduli stacks/spaces
%
\def\QCohstack{\mathcal{QC}\!\mathit{oh}}
\def\Cohstack{\mathcal{C}\!\mathit{oh}}
\def\Spacesstack{\mathcal{S}\!\mathit{paces}}
\def\Quotfunctor{\mathrm{Quot}}
\def\Hilbfunctor{\mathrm{Hilb}}
\def\Curvesstack{\mathcal{C}\!\mathit{urves}}
\def\Polarizedstack{\mathcal{P}\!\mathit{olarized}}
\def\Complexesstack{\mathcal{C}\!\mathit{omplexes}}
% \Pic is the operator that assigns to X its picard group, usage \Pic(X)
% \Picardstack_{X/B} denotes the Picard stack of X over B
% \Picardfunctor_{X/B} denotes the Picard functor of X over B
\def\Pic{\mathop{\mathrm{Pic}}\nolimits}
\def\Picardstack{\mathcal{P}\!\mathit{ic}}
\def\Picardfunctor{\mathrm{Pic}}
\def\Deformationcategory{\mathcal{D}\!\mathit{ef}}
\usepackage{cleveref}
\crefname{originaltheorem}{original theorem}{original theorems}
\Crefname{originaltheorem}{Original theorem}{Original theorems}
\makeatletter
\newlabel{theorem-injective-embedding-grothendieck}{{079H}{source}{}{sourcepointer-1}{}}
\newlabel{theorem-injective-embedding-grothendieck@cref}{{[originaltheorem][0][]079H}{[1][0][]source}{}{}{}}
\makeatother
\title{Functorial K-injective replacements with injective terms}
\author{The Stacks Project Authors}
\date{stacks-project:injectives.tex:local-snapshot-20261012; GNU FDL 1.2 or later}
\begin{document}
\maketitle
\noindent\textit{Editorial note.} All delivered mathematical text is the author's own, copied from the pinned Stacks Project chapter file sites-cohomology.tex; no statement or proof step was rewritten, reordered or removed.\par\medskip
\noindent\textit{Editorial note.} The proof is the author's complete construction, each of the two asserted properties being verified in the span; the closing sentence explains how the quasi-isomorphism is checked and is preserved verbatim.\par\medskip
\noindent\textit{Editorial note.} Every labelled reference occurring inside the counted spans is carried in assembly.external\_references as a non-load-bearing navigation pointer (same-chapter ones at their definition site, cross-chapter ones at the invocation inside the counted span), so no load-bearing mathematics is replaced by a stub. No bibliography citation occurs inside the counted spans.\par\medskip
\noindent\textit{Editorial note.} The same-name check for this card scanned every staged Stacks chapter of this repository -- all sources/stacks-*/pinned-original-source.tex files, the card's own chapter excluded -- and the labels found are recorded in dedup.same\_name\_cross\_chapter on the card.\par\medskip
\noindent\textit{Editorial note.} The preamble is the chapter's own retained minimal preamble (sources/stacks-sites-cohomology/retained-preamble.tex), a verbatim excerpt of the native preamble: the Xy-pic setup and the native theorem declarations with the mathematical macros.\par\medskip
\noindent\textit{Allowed standard prerequisites (editorial).}\begin{itemize}
\item Functorial injective embeddings. Same chapter, Theorem theorem-injective-embedding-grothendieck. Supplies the functorial embedding of an object into an injective object that the construction starts from.
\item The two-term complex attached to an object. Same chapter, the construction of the functor M from objects to complexes and the description of its terms. Supplies the functorial complexes and the short exact sequences of complexes used to compare the given complex with its replacement.
\item Products of injectives and the long exact cohomology sequence. Homology, Lemma homology-lemma-product-injectives and the long exact cohomology sequence. Supply that products of injective objects are injective and that the comparison maps induce isomorphisms on cohomology.
\end{itemize}\medskip
\ifdefined\hypertarget\hypertarget{sourcepointer-1}{}\fi
\noindent\textit{Source reference.} Original theorem 079H: sources/stacks-injectives/pinned-original-source.tex, source lines 1716--1716. Navigation pointer invoked at proof line 1976; the counted span contains the pointer only and the cited result is not reproduced inside it, so the pointer is registered as a non-load-bearing navigation stub in the established form for Stacks material. No bibliography citation is involved.\par\medskip
\begin{lemma}
\label{lemma-functorial-injective}
Let $\mathcal{A}$ be a Grothendieck abelian category.
There exists a functor $M^\bullet \mapsto \mathbf{N}^\bullet(M^\bullet)$
and a natural transformation
$j_{M^\bullet} : M^\bullet \to \mathbf{N}^\bullet(M^\bullet)$
such
\begin{enumerate}
\item $j_{M^\bullet}$ is a (termwise) injective quasi-isomorphism, and
\item for every $n \in \mathbf{Z}$ the map $M^n \to \mathbf{N}^n(M^\bullet)$
factors through a subobject $I^n \subset \mathbf{N}^n(M^\bullet)$ where $I^n$
is an injective object of $\mathcal{A}$.
\end{enumerate}
\end{lemma}

\begin{proof}
Choose a functorial injective embeddings $i_M : M \to I(M)$, see
Theorem \ref{theorem-injective-embedding-grothendieck}.
For every complex $M^\bullet$ denote $J^\bullet(M^\bullet)$ the complex
with terms $J^n(M^\bullet) = I(M^n) \oplus I(M^{n + 1})$ and differential
$$
d_{J^\bullet(M^\bullet)} =
\left(
\begin{matrix}
0 & 1 \\
0 & 0
\end{matrix}
\right)
$$
There exists a canonical injective map of complexes
$u_{M^\bullet} : M^\bullet \to J^\bullet(M^\bullet)$ by mapping $M^n$ to
$I(M^n) \oplus I(M^{n + 1})$ via the maps $i_{M^n} : M^n \to I(M^n)$ and
$i_{M^{n + 1}} \circ d : M^n \to M^{n + 1} \to I(M^{n + 1})$. Hence a
short exact sequence of complexes
$$
0 \to M^\bullet \xrightarrow{u_{M^\bullet}}
J^\bullet(M^\bullet) \xrightarrow{v_{M^\bullet}}
Q^\bullet(M^\bullet) \to 0
$$
functorial in $M^\bullet$. Set
$$
\mathbf{N}^\bullet(M^\bullet) = C(v_{M^\bullet})^\bullet[-1].
$$
Note that
$$
\mathbf{N}^n(M^\bullet) = Q^{n - 1}(M^\bullet) \oplus J^n(M^\bullet)
$$
with differential
$$
\left(
\begin{matrix}
- d^{n - 1}_{Q^\bullet(M^\bullet)} & - v^n_{M^\bullet} \\
0 & d^n_{J^\bullet(M)}
\end{matrix}
\right)
$$
Hence we see that there is a map of complexes
$j_{M^\bullet} : M^\bullet \to \mathbf{N}^\bullet(M^\bullet)$
induced by $u$. It is injective and factors through an injective subobject
by construction. The map $j_{M^\bullet}$ is a quasi-isomorphism as one
can prove by looking at the long exact sequence of cohomology associated
to the short exact sequences of complexes above.
\end{proof}

\end{document}
